% | Theory consequence            | Metric |
% | ----------------------------- | ------ |
% | local compactness             | Cosine |
% | relational geometry stability | CKA    |
% | dominant semantic directions  | PWCCA  |
% | compact low-dim region        | IFAC   |
\subsection{Empirical Evaluation of Invariant Zones}
\label{sec: metric}
The geometric formulation of invariant latent features in Section~\ref{sec: define} yields several observable consequences for representations produced by a model under semantic-preserving transformations. If paraphrasing primarily induces motion along nuisance directions of the tangent space while preserving semantic identity, then latent representations should exhibit consistent geometric structure across paraphrase variants. In particular, four measurable properties follow from the local geometric hypothesis.

First, paraphrases expressing the same meaning should remain locally compact in representation space, reflecting stability of the underlying semantic point on the manifold. Second, although individual representations may move within the tangent neighborhood, the relational geometry among inputs should remain stable, preserving the organization of semantic clusters. Third, the dominant latent directions that encode semantic information should remain aligned across paraphrase variants. Finally, paraphrases corresponding to the same meaning should occupy a compact, low-dimensional region of the representation manifold.

To operationalize these predictions, we construct a controlled paraphrase evaluation setting. For each semantic meaning class, we collect multiple paraphrase variants that preserve semantic intent while varying lexical form and syntactic structure. Let $\mathcal{W}$ denote the set of all within-class paraphrase pairs and $\mathcal{A}$ denote the set of across-class pairs drawn from distinct meaning classes. This within–across formulation allows us to distinguish semantic stability from trivial representational collapse: invariant representations must remain stable within meaning classes while preserving separation between different meanings.

We then evaluate invariance using four complementary metrics that probe different geometric aspects of the representation space. Each metric corresponds to one of the theoretical consequences derived from the local geometric hypothesis. Cosine similarity measures local semantic compactness of paraphrase representations. Centered Kernel Alignment (CKA) evaluates the stability of relational geometry across transformation variants. Projection-Weighted Canonical Correlation Analysis (PWCCA) measures the alignment of dominant latent subspaces that encode semantic information. Finally, the Intrinsic Feature Alignment Coefficient (IFAC) quantifies the intrinsic dimensionality and compactness of paraphrase clusters. A summary is presented in table~\ref{tab:metrics_summary}. 

\begin{table}[ht]
\centering
\vspace{-0.5em}
\begin{tabularx}{0.95\linewidth}{p{0.68\linewidth} X}
\toprule
\textbf{Geometric consequence of invariance hypothesis} & \textbf{Metric probe} \\
\midrule
Paraphrases of the same meaning remain locally compact in representation space. 
& Cosine Similarity \\

Relational geometry among inputs remains stable under semantic-preserving transformations. 
& CKA \\

Dominant latent directions encoding semantic information remain aligned across paraphrase variants. 
& PWCCA \\

Paraphrases corresponding to the same meaning occupy a compact, low-dimensional region of the representation manifold. 
& IFAC \\
\bottomrule
\end{tabularx}
\vspace{0.5em}
\caption{Connection between the geometric consequences of the invariant latent feature hypothesis and the empirical metrics used to probe them. Each metric evaluates a different aspect of the predicted representation geometry.}
\label{tab:metrics_summary}
\vspace{-1.2em}
\end{table}

Together, these metrics provide a multi-perspective empirical probe of the geometric invariance hypothesis. By analyzing their behavior across model layers, we can identify regions where semantic representations become stable under paraphrasing while remaining discriminative across meanings. These regions correspond to the invariant zones hypothesized in Section~\ref{sec: define} and form the basis for the layer-wise analyses and attribution experiments presented in subsequent sections.

\subsubsection{Cosine Similarity}

We first quantify invariance using cosine similarity, which measures the local compactness of paraphrase representations in latent space. For a given layer $\ell$, cosine similarity evaluates whether representations corresponding to semantically equivalent inputs cluster more tightly than those corresponding to semantically distinct inputs.

Let $\mathcal{W}$ denote the set of all paraphrase pairs drawn from the same meaning class, and let $\mathcal{A}$ denote the set of all pairs drawn from different meaning classes. We define
\[
\text{cosine}_{\text{within}} = \mathbb{E}_{(x_i,x_j)\in \mathcal{W}}
\left[ 
\frac{\langle \Phi_\ell(x_i), \Phi_\ell(x_j)\rangle}
{\|\Phi_\ell(x_i)\|\;\|\Phi_\ell(x_j)\|}
\right],
\]
and
\[
\text{cosine}_{\text{across}} = \mathbb{E}_{(x_i,x_j)\in \mathcal{A}}
\left[ 
\frac{\langle \Phi_\ell(x_i), \Phi_\ell(x_j)\rangle}
{\|\Phi_\ell(x_i)\|\;\|\Phi_\ell(x_j)\|}
\right].
\]

The cosine-based invariance index is then defined as
\begin{equation}
\text{cosine}_{\text{InvIdx}} =
\frac{\text{cosine}_{\text{within}} - \text{cosine}_{\text{across}}}
{\sigma(\text{cosine}_{\text{across}}) + \epsilon},
\end{equation}
where $\sigma(\text{cosine}_{\text{across}})$ denotes the standard deviation of across-class similarities and $\epsilon$ is a small constant for numerical stability.

A high $\text{cosine}_{\text{InvIdx}}$ indicates that paraphrased inputs occupy a compact region in latent space while remaining well separated from representations of different meanings. This behavior is consistent with the geometric hypothesis that semantic-preserving transformations primarily induce motion along nuisance directions that remain within the same semantic neighborhood on the representation manifold.

\subsubsection{Centered Kernel Alignment}

Next, we characterize whether the internal geometric organization of latent representations is preserved under semantic-preserving transformations. 
To analyze invariance at the level of relational geometry, we employ Centered Kernel Alignment (CKA), which measures structural similarity between two sets of representations.

Let $X, Y \in \mathbb{R}^{N \times d}$ denote activation matrices constructed from $N$ inputs, where each row corresponds to a pooled latent representation $\Phi_\ell(x)$ at a fixed layer $\ell$, obtained from two transformation variants of the same input set (e.g., original prompts and their paraphrases). We define the associated linear kernel matrices as
\(
K_X = X X^\top, \quad K_Y = Y Y^\top.
\)
Each kernel is centered using the centering matrix $H = I - \frac{1}{N}\mathbf{1}\mathbf{1}^\top$, yielding centered kernels $HK_XH$ and $HK_YH$. The linear CKA similarity is then computed as
\begin{equation}
\text{CKA}(X, Y) =
\frac{\langle HK_XH, HK_YH \rangle_F}
{\|HK_XH\|_F \, \|HK_YH\|_F},
\end{equation}
where $\langle \cdot, \cdot \rangle_F$ denotes the Frobenius inner product.

A high CKA value indicates that pairwise relationships among inputs are preserved across paraphrase variants. In the context of the geometric invariance hypothesis, this implies that semantic-preserving transformations maintain the relational organization of meanings on the representation manifold even when individual representations undergo local perturbations.

\subsubsection{Projection-Weighted Canonical Correlation Analysis}

Next, we explicitly evaluates whether the dominant latent directions that encode semantic information remain stable across transformations. 
To analyze invariance at the level of informative subspaces, we employ Projection-Weighted Canonical Correlation Analysis (PWCCA).

Let $X, Y \in \mathbb{R}^{N \times d}$ denote mean-centered activation matrices at layer $\ell$, corresponding to two transformation variants of the same input set. Canonical Correlation Analysis identifies pairs of projection vectors $(w_i, v_i)$ that maximize the correlation between projected components,
\(
\rho_i = \mathrm{corr}(w_i^\top X, v_i^\top Y), \quad i = 1, \ldots, k,
\)
where $k$ denotes the number of canonical components.

PWCCA aggregates these canonical correlations by weighting each component according to its contribution to the variance of the representation,
\begin{equation}
\text{PWCCA}(X, Y) =
\sum_{i=1}^{k} \alpha_i \rho_i, \quad
\alpha_i =
\frac{\mathbb{E}[|X w_i|]}{\sum_j \mathbb{E}[|X w_j|]}.
\end{equation}

High PWCCA values indicate that the dominant latent subspaces of the two representation sets remain aligned. Under the geometric invariance hypothesis, this suggests that semantic-preserving transformations maintain the principal directions along which semantic information is encoded, while variation induced by paraphrasing occurs primarily in secondary nuisance directions.

\subsubsection{Intrinsic Feature Alignment Coefficient}

Finally, we characterize the internal geometric structure of paraphrase clusters themselves. To quantify the intrinsic compactness of such clusters, we introduce the Intrinsic Feature Alignment Coefficient (IFAC), which measures the extent to which variance within a paraphrase group is concentrated in a low-dimensional subspace.

For each meaning group $g$, let $X_g \in \mathbb{R}^{N_g \times d}$ denote the matrix of pooled latent representations corresponding to paraphrases of the same semantic meaning. We perform principal component analysis on the mean-centered $X_g$ and obtain singular values $\{\sigma_i\}$. The explained variance ratio of the $i$-th principal component is given by
% \(
% \mathrm{EVR}_i = \frac{\sigma_i^2}{\sum_j \sigma_j^2}.
% \)
\(
\mathrm{EVR}_i = \sigma_i^2 / \sum_j \sigma_j^2.
\)

The IFAC score for group $g$ at layer $\ell$ is defined as the cumulative explained variance captured by the top-$k$ principal components,
\begin{equation}
\text{IFAC}_\ell^{(g)} =
\sum_{i=1}^{k} \mathrm{EVR}_i,
\end{equation}
and the layer-wise IFAC is obtained by averaging across all meaning groups,
\begin{equation}
\text{IFAC}_\ell =
\mathbb{E}_g\!\left[\text{IFAC}_\ell^{(g)}\right].
\end{equation}

A high IFAC value indicates that paraphrases of the same meaning lie within a compact, low-dimensional region of the representation manifold. This behavior is consistent with the geometric hypothesis that semantic-preserving transformations move representations along nuisance directions within a localized neighborhood while leaving the semantic identity of the representation largely unchanged.